\section{Дифференциальные уравнения первого порядка}

\begin{definition}
\emph{ОДУ первого порядка} — уравнение $y'=f(x,y)$. \emph{Общее решение} — семейство
$y=\varphi(x,C)$, зависящее от произвольной постоянной; \emph{частное} получается фиксацией $C$ из
\emph{начального условия} $y(x_0)=y_0$ (\emph{задача Коши}); \emph{особое} решение не входит в общее
семейство ни при каком $C$ (огибающая).
\end{definition}

\begin{theorem}[существование и единственность]
Если $f(x,y)$ непрерывна и удовлетворяет условию Липшица по $y$
($\abs{f(x,y_1)-f(x,y_2)}\le L\abs{y_1-y_2}$) в окрестности $(x_0,y_0)$, то задача Коши имеет
единственное решение в некоторой окрестности $x_0$.
\end{theorem}

\begin{remark}[разделяющиеся переменные]
Уравнение $y'=g(x)h(y)$ решается разделением: $\dfrac{dy}{h(y)}=g(x)\dx$, $\int\dfrac{dy}{h(y)}=\int g(x)\dx+C$.
Условие Липшица по $y$ (например, гладкость $h$) гарантирует единственность; в нулях $h$ возможны
особые решения.
\end{remark}

\section{Уравнения, допускающие понижение порядка}

\begin{property}[нет явной $y$: $y''=f(x,y')$]
Подстановка $p=y'$ понижает порядок: $p'=f(x,p)$ — уравнение первого порядка относительно $p(x)$;
найдя $p$, интегрируем $y=\int p\,dx$.
\end{property}

\begin{property}[нет явной $x$: $y''=f(y,y')$]
Берём $p=y'$ как функцию от $y$: $y''=\dfrac{dp}{dx}=\dfrac{dp}{dy}\,p$. Получаем уравнение первого
порядка $p\dfrac{dp}{dy}=f(y,p)$ относительно $p(y)$.
\end{property}

\section{Линейная зависимость. Вронскиан}

\begin{definition}
Функции $y_1,\dots,y_n$ \emph{линейно независимы} на интервале, если
$\sum c_k y_k(x)\equiv0\Rightarrow$ все $c_k=0$. \emph{Определитель Вронского}
\[
  W(x)=\begin{vmatrix} y_1 & \dots & y_n\\ y_1' & \dots & y_n'\\ \vdots & & \vdots\\
  y_1^{(n-1)} & \dots & y_n^{(n-1)}\end{vmatrix}.
\]
\end{definition}

\begin{remark}
Если $W(x_0)\ne0$ хотя бы в одной точке, функции линейно независимы (обратное для произвольных
функций неверно; для решений ЛОУ — верно).
\end{remark}

\begin{theorem}[формула Остроградского–Лиувилля]
Для решений уравнения $y''+p(x)y'+q(x)y=0$ вронскиан удовлетворяет $W'=-p(x)W$, откуда
\[
  W(x)=W(x_0)\exp\Bigl(-\int_{x_0}^x p(t)\,dt\Bigr).
\]
\end{theorem}

\begin{proof}
$W=y_1y_2'-y_2y_1'$, поэтому $W'=y_1y_2''-y_2y_1''$ (члены $y_1'y_2'$ сокращаются). Подставив
$y_i''=-p y_i'-q y_i$:
\[
  W'=y_1(-p y_2'-q y_2)-y_2(-p y_1'-q y_1)=-p(y_1y_2'-y_2y_1')=-pW .
\]
Интегрируя линейное уравнение $W'=-pW$, получаем экспоненциальную формулу.
\end{proof}

\begin{corollary}
$W$ либо тождественно равен нулю, либо нигде не обращается в нуль. \emph{Критерий ФСР:} решения
$y_1,\dots,y_n$ ЛОУ образуют фундаментальную систему $\iff W\ne0$ (тогда общее решение
$y=\sum C_k y_k$).
\end{corollary}

\section{ЛОУ с постоянными коэффициентами}

Рассмотрим $y^{(n)}+a_1y^{(n-1)}+\dots+a_n y=0$.

\begin{theorem}[метод Эйлера]
Подстановка $y=e^{\lambda x}$ приводит к \emph{характеристическому уравнению}
\[
  \lambda^n+a_1\lambda^{n-1}+\dots+a_n=0 .
\]
Каждому простому вещественному корню $\lambda$ отвечает решение $e^{\lambda x}$; паре
комплексно-сопряжённых $\alpha\pm i\beta$ — решения $e^{\alpha x}\cos\beta x$,
$e^{\alpha x}\sin\beta x$; корню кратности $m$ — решения $e^{\lambda x},xe^{\lambda x},\dots,x^{m-1}e^{\lambda x}$.
\end{theorem}

\begin{proof}
$\frac{d}{dx}e^{\lambda x}=\lambda e^{\lambda x}$, поэтому подстановка $y=e^{\lambda x}$ даёт
$e^{\lambda x}\,p(\lambda)=0$, где $p$ — характеристический многочлен; значит $e^{\lambda x}$ —
решение $\iff p(\lambda)=0$. Для комплексного корня берут вещественную и мнимую части.
\end{proof}

\begin{theorem}[ФСР при кратных корнях]
Функции $e^{\lambda x},xe^{\lambda x},\dots,x^{m-1}e^{\lambda x}$ линейно независимы.
\end{theorem}

\begin{proof}
Если $\sum_{k=0}^{m-1}c_k x^k e^{\lambda x}\equiv0$, поделим на $e^{\lambda x}\ne0$:
$\sum c_k x^k\equiv0$ — многочлен, тождественно равный нулю, поэтому все $c_k=0$. (Что эти функции —
решения при корне кратности $m$, проверяется подстановкой: оператор $\bigl(\tfrac{d}{dx}-\lambda\bigr)^m$
обнуляет $x^k e^{\lambda x}$ при $k<m$.)
\end{proof}

\begin{remark}[физическая интерпретация]
Для $y''+2\delta y'+\omega_0^2 y=0$: корни $\lambda=-\delta\pm\sqrt{\delta^2-\omega_0^2}$. При
$\delta>\omega_0$ — два вещественных корня (\emph{апериодический} режим), при $\delta=\omega_0$ —
кратный корень (\emph{критический}), при $\delta<\omega_0$ — комплексные корни
(\emph{колебательный} режим, затухающие колебания).
\end{remark}

\section{Нормальные системы ОДУ}

\begin{definition}
\emph{Нормальная система}: $\dot{\mathbf x}=F(t,\mathbf x)$, $\mathbf x\in\R^n$. \emph{Фазовое
пространство} — пространство значений $\mathbf x$; \emph{фазовая траектория} — кривая решения в нём;
\emph{фазовый портрет} — совокупность траекторий.
\end{definition}

\begin{remark}[метод исключения]
Систему $2$-го порядка сводят к одному уравнению: из первого уравнения выражают одну переменную,
дифференцируют и подставляют, получая ЛОУ $2$-го порядка относительно другой.
\end{remark}

\begin{theorem}[матричный метод Эйлера]
Для линейной системы $\dot{\mathbf x}=A\mathbf x$ ищем решения вида $\mathbf x=e^{\lambda t}\mathbf v$.
Подстановка даёт $A\mathbf v=\lambda\mathbf v$ — характеристическое уравнение $\det(A-\lambda I)=0$.
\end{theorem}

\begin{remark}
ФСР строится по собственным значениям $A$: вещественным различным — $e^{\lambda_k t}\mathbf v_k$;
комплексным $\alpha\pm i\beta$ — вещественная и мнимая части $e^{(\alpha+i\beta)t}\mathbf v$;
кратным — с добавлением множителей $t,\dots,t^{k-1}$ (присоединённые векторы).
\end{remark}

\section{Теория устойчивости}

\begin{definition}[по Ляпунову]
Положение равновесия $\mathbf x_0$ (где $F(\mathbf x_0)=0$) \emph{устойчиво}, если
$\forall\eps>0\ \exists\delta>0$: из $\norm{\mathbf x(0)-\mathbf x_0}<\delta$ следует
$\norm{\mathbf x(t)-\mathbf x_0}<\eps$ при всех $t\ge0$. \emph{Асимптотически устойчиво}, если вдобавок
$\mathbf x(t)\to\mathbf x_0$. \emph{Неустойчиво} — если условие устойчивости нарушено.
\end{definition}

\begin{theorem}[устойчивость по первому приближению]
Пусть $A=F'(\mathbf x_0)$ — матрица линеаризации. Если все собственные значения имеют отрицательную
вещественную часть ($\Re\lambda_i<0$), то равновесие асимптотически устойчиво; если хотя бы у одного
$\Re\lambda_i>0$ — неустойчиво.
\end{theorem}

\begin{remark}[топологическая классификация на плоскости]
Для $\dot{\mathbf x}=A\mathbf x$ в $\R^2$ по собственным значениям $\lambda_{1,2}$:
\begin{itemize}[leftmargin=*,nosep]
  \item вещественные одного знака — \emph{узел} (устойчивый при $\lambda<0$, неустойчивый при $\lambda>0$);
  \item вещественные разных знаков — \emph{седло} (всегда неустойчиво);
  \item комплексные с $\Re\lambda\ne0$ — \emph{фокус} (закручивающаяся спираль; устойчив при $\Re\lambda<0$);
  \item чисто мнимые $\pm i\beta$ — \emph{центр} (замкнутые траектории, устойчив, но не асимптотически).
\end{itemize}
\end{remark}
